\documentclass[11pt]{article}
\usepackage[margin=0.85in]{geometry}
\usepackage{amsmath,amssymb,booktabs,tabularx,array,xcolor,hyperref}
\usepackage[T1]{fontenc}
\usepackage{microtype}
\hypersetup{colorlinks=true,linkcolor=blue,citecolor=blue,urlcolor=blue}
\newcommand{\system}{ActiveRetina-RC}
\newcommand{\statusnote}{\par\smallskip\noindent\fbox{\parbox{0.96\linewidth}{\textbf{Evidence status.} This is a retrospective development study. The 70.13\% result is a same-cohort, score-level replay with thresholds and rank-reference distributions derived from an inspected cohort. It is not independent validation, a clinical certificate, or an image-to-diagnosis product evaluation.}}\par\smallskip}
\title{Retrospective Zero-Miss Specificity in Multi-Source Diabetic-Retinopathy Screening: Replay, Patient-Grouped Recalibration, and Domain-Transfer Stress Tests}
\author{Burak Aksu}
\date{October 2026}
\begin{document}
\maketitle
\begin{abstract}
Retinal screening studies often emphasize sensitivity or area under the receiver-operating-characteristic curve while leaving the operating-point selection process difficult to audit. We report an evidence audit of a frozen, multi-source diabetic-retinopathy score policy developed from public DDR, DeepDRiD, and DRTiD data. On a repeatedly inspected pool of 981 eye-level cases, the policy assigns a binary CLEAR or REFER action to every supplied score record, clears 270 of 385 gradable nonreferable eyes (70.13\% specificity), and produces no observed false-negative clearances among 463 gradable referable eyes; all 133 ungradable cases are referred. This is a same-cohort replay: source-specific percentile-reference arrays and zero-miss cutoffs were derived from this pool. Five-fold patient-grouped recalibration yields 268/385 true negatives (69.61\%) and three referable-or-ungradable misses; nested weight/threshold selection yields 248/385 (64.42\%) with five misses, while a nested margin variant yields 253/385 (65.71\%) with six. A post-hoc zero-miss frontier reaches 243/385 (63.12\%), but its margins were selected on held-out predictions. Separate MMRDR-UWF stress tests reach 34.50\% specificity at a same-validation-set zero-miss point, and an MMRDR-initialized fine-tuning run is exposed to the reference cohort through a corrected manifest overlap. Neither is independent confirmation. A transfer candidate OR-combined with the frozen policy adds three clearances but introduces two misses. The contribution is a reproducible audit separating deterministic replay, internal development estimation, and independent external validation. We make no independent clinical-performance, risk-certification, or complete image-to-diagnosis claim.
\end{abstract}

\noindent\textbf{Keywords:} diabetic retinopathy, fundus photography, specificity, patient-level validation, score fusion, selective classification, data leakage, risk control.

\statusnote

\section{Introduction}
Diabetic retinopathy (DR) screening is a thresholded decision problem, not merely a ranking problem. A model may rank referable disease well yet operate poorly when a system must decide which eyes to clear and which to refer. This distinction becomes especially important when investigators seek a zero-observed-miss operating point: each threshold adjustment trades missed disease against false referrals, and repeated inspection of the evaluation cohort can make that trade appear more reliable than it is. The broader methodological value is to quantify how a reported operating point changes across three distinct evidence classes—deterministic replay, internal development estimation, and independent external validation—so that repeated policy selection in any multi-source medical pipeline is not mistaken for transportable performance.

The present study audits a concrete development result in \system: a pooled score-level policy across three public fundus datasets. Its attractive headline is 70.13\% specificity, zero observed referable misses, and 100\% binary action assignment. Those values are mathematically reproducible. The scientific question is whether they survive when source-specific thresholds and blend weights are estimated without using the evaluated patients. The answer is materially different: patient-grouped procedures reduce specificity and introduce missed referable or ungradable cases.

We therefore frame this work as a reproducibility and validation-methods report rather than a new diagnostic algorithm. The principal contribution is the traceable comparison between (i) deterministic replay of a same-cohort-selected policy and (ii) patient-grouped recalibration and nested selection on the same previously inspected cohort. A further transfer diagnostic tests whether a separately trained candidate improves the frozen baseline without sacrificing its zero-miss behavior.

\subsection{Contributions}
\begin{enumerate}
\item We specify and reproduce the frozen source-conditional rank-fusion rule, including feature families, weights, thresholds, and the DDR quality override.
\item We report the exact composition of the 981-case replay and source-specific specificity counts, keeping gradable referable, gradable nonreferable, and ungradable labels distinct.
\item We compare the headline replay with five-fold patient-grouped calibration, nested patient-fold weight/threshold selection, and a clearly identified post-hoc zero-miss frontier.
\item We document a later clean-development-trained transfer diagnostic that gains three true negatives under an OR rule but incurs two misses, illustrating why a point-estimate lift is not sufficient.
\item We report a distinct MMRDR-UWF transfer stress test and a provenance correction for a multi-source fine-tuning run whose validation manifest overlapped the historical reference cohort.
\item We provide a fail-closed interpretation and reproducibility map that prohibit treating replay fidelity or exposed-cohort results as clinical certification.
\end{enumerate}

\section{Related Work}
Deep-learning systems have achieved strong DR-screening performance in large retinal-photograph studies, including multiethnic and primary-care evaluations \cite{gulshan2016,ting2017,abramoff2018,bellemo2019}. Public datasets such as DDR and DeepDRiD support reproducible work on grading, image quality, and cross-domain behavior \cite{ddr,deepdrid}. These works motivate rigorous separation between model discrimination, operating-point selection, and evaluation. Model-selection overfitting and selection bias on finite evaluation sets are established concerns \cite{cawley2010,varma2006}; patient-grouped folds are important when observations share a patient hierarchy \cite{roberts2017}.

Selective prediction formalizes the option to abstain when a prediction is not sufficiently reliable \cite{geifman2017,selectivenet}. Risk-control methods can provide finite-sample guarantees under prespecified procedures and assumptions; empirical threshold search alone does not do so \cite{vovk2005,angelopoulos2021,angelopoulos2024}. Prediction-model reporting guidance emphasizes transparent data partitioning, evaluation, and limitations \cite{tripodai,claim,claim2024}. The present audit applies those principles to a pooled specificity figure whose thresholds were selected on the same cohort used for reporting.

\section{Study Question and Endpoints}
For eye-level case $i$, let $Y_i\in\{0,1\}$ denote gradable referable DR, with $Y_i=1$ corresponding to referable disease (dataset grades 2--4 under the harmonized project definition), and let $U_i=1$ mark an ungradable case. The policy emits $A_i\in\{\text{CLEAR},\text{REFER}\}$. We report specificity among gradable nonreferable cases,
\begin{equation}
\mathrm{Sp}=\frac{\sum_i\mathbb{1}(Y_i=0,U_i=0,A_i=\text{CLEAR})}{\sum_i\mathbb{1}(Y_i=0,U_i=0)},
\end{equation}
missed referable cases $M=\sum_i\mathbb{1}(Y_i=1,U_i=0,A_i=\text{CLEAR})$, missed ungradable cases $M_U=\sum_i\mathbb{1}(U_i=1,A_i=\text{CLEAR})$, and action assignment $C=\frac{1}{N}\sum_i\mathbb{1}(A_i\in\{\text{CLEAR},\text{REFER}\})$. Here $C=100\%$ means every record receives one of the two computational actions; it does not imply independent generalization, image-quality adequacy, or clinical usability.

The prespecified development target used in the project was a useful specificity at zero observed misses and full binary action assignment. This paper does not present that target as a validated safety guarantee. We use ``zero-miss'' only as a finite-sample descriptive observation unless explicitly qualified otherwise.

\section{Data and Provenance}
\subsection{Inspected pooled reference cohort}
The frozen replay contains 981 eye-level records: 625 DDR, 106 DeepDRiD, and 250 DRTiD. The label composition is 385 gradable nonreferable, 463 gradable referable, and 133 ungradable cases. The cases and score outputs were repeatedly inspected during model and threshold development. The cohort is therefore exposed and cannot serve as an untouched final test.

\begin{table}[t]
\centering
\caption{Frozen replay cohort and source-specific results. All rows received a binary action; specificity denominators include gradable nonreferable cases only.}
\label{tab:cohort}
\begin{tabular}{lrrrrrrr}
\toprule
Source & Cases & Nonref. & Ref. & Ungrad. & TN & FP & Sp. \\
\midrule
DDR & 625 & 206 & 286 & 133 & 134 & 72 & 65.05\% \\
DeepDRiD & 106 & 52 & 54 & 0 & 51 & 1 & 98.08\% \\
DRTiD & 250 & 127 & 123 & 0 & 85 & 42 & 66.93\% \\
\midrule
Pooled & 981 & 385 & 463 & 133 & 270 & 115 & 70.13\% \\
\bottomrule
\end{tabular}
\end{table}

The datasets have different acquisition and label characteristics \cite{ddr,deepdrid,drtidpaper}; domain generalization remains an active challenge in DR classification \cite{galappaththige2024}. Pooling gives each eye one vote in the pooled denominator; it is not a population-standardized estimate. Dataset availability and use remain subject to the corresponding licenses and access conditions. We do not redistribute raw images.

\subsection{Exposure and split history}
The score references and operating thresholds in the frozen policy were derived from the same repeatedly inspected 981-case cohort. Thus, the replay answers ``can the recorded recipe reproduce its recorded result?'' It does not answer ``what performance should be expected on new patients?'' Patient-grouped calibration folds below prevent fold-specific calibration from using the held-out patients, but cannot erase prior inspection of the cohort or prior global model-family exploration. We label all such results development diagnostics.

\section{Frozen Policy}
For source $d$, each component score $z_{idm}$ is transformed into a percentile midrank $r_{idm}$ against the stored, sorted source-specific reference vector of size $n_d$:
\begin{equation}
r_{idm}=\frac{L(z_{idm})+R(z_{idm})+1}{2n_d},
\end{equation}
where $L$ and $R$ count reference values strictly less than and less than-or-equal to the score. The rank-fused risk score is a source-specific convex combination:
\begin{align}
s_{\mathrm{DDR}}&=0.50r_{\mathrm{APTOS}}+0.10r_{\mathrm{ordinal}}+0.40r_{\mathrm{audit}},\\
s_{\mathrm{DeepDRiD}}&=0.20r_{\mathrm{ordinal-quality}}+0.40r_{\mathrm{source-balanced}}+0.40r_{\mathrm{fundus}},\\
s_{\mathrm{DRTiD}}&=0.72r_{\mathrm{rank-pair}}+0.08r_{\mathrm{fundus-margin}}+0.20r_{\mathrm{geometry}}.
\end{align}
An eye is referred when its score reaches its source-specific cut point; the DDR rule also refers cases whose quality score reaches its stored override threshold. The cut points are 0.22224 for DDR, 0.435849 for DeepDRiD, and 0.34496 for DRTiD; the DDR quality override is 0.0018599. These numerical cut points and rank-reference vectors come from the exposed cohort. They are published here to make the replay auditable, not as proposed portable clinical thresholds.

The implementation is a deterministic score-policy evaluator. It requires three precomputed component outputs per source and a source identifier. It is not a single end-to-end image model. Exact source checkpoints and preprocessing implementations are not all present in the current repository state; accordingly, the replay should not be described as complete raw-image inference. Throughout, “deterministic replay” means recomputing actions from supplied component-score rows and a frozen policy; “internal development estimation” means grouped or nested analysis on the previously inspected project cohorts; and “independent external validation” means evaluation on a new cohort untouched by training, calibration, threshold/policy selection, or prior project decisions. No independent external validation is reported here.

\section{Evaluation Procedures}
\subsection{Frozen same-cohort replay}
We reran the frozen policy on each of the 981 stored component-score rows and compared every emitted action and risk score with the recorded replay. We then reconstructed source-level confusion counts and checked the headline grouped/nested counts against their saved summary artifacts. The executable command, required file inventory, environment dependency, checksums, and scope are specified in the accompanying reproduction-package README. This validates deterministic replay of supplied scores and manuscript transcription; it does not regenerate image features, retrain component models, or provide independent external validation.

\subsection{Patient-grouped calibration and nested selection}
The fixed-recipe analysis used five patient-grouped folds. For each fold, source-specific calibration was performed without its held-out patients, and predictions were aggregated over the held-out folds. The component model family and the cohort had already been inspected globally; this is not an independent confirmatory evaluation.

The nested analysis used five outer patient-grouped folds and four inner folds. Blend weights were searched on a simplex grid with 0.1 increments; the risk threshold was the minimum required-positive score in the training portion. The inner search compared weight/margin combinations, prioritizing zero inner misses before specificity; tested margins were 0, 0.005, 0.01, 0.02, 0.03, 0.05, 0.075, 0.1, 0.15, and 0.2. The selected combination was evaluated on its outer fold. This protects each outer fold from that run's direct fitting, but the data and component family remain previously exposed. Separately, a post-hoc margin sweep searched 13 values from 0 to 0.2 and chose source-specific margins after inspecting held-out predictions; that frontier is included solely to show the attainable retrospective trade-off and is not an unbiased candidate.

\subsection{Clean-development candidate transfer diagnostic}
A ResNet-50 initialized from a public MMRDR-UWF model was fine-tuned at 448-pixel input on original training partitions after excluding all patient groups mapped to the 981 reference cases. Its patient-disjoint development split selected the checkpoint and two thresholds. The frozen candidate was then evaluated on the repeatedly inspected reference cohort. We also computed an offline union action, CLEAR if either the old baseline or candidate cleared. Because the reference cohort has extensive prior project exposure, this is a transfer diagnostic—not an independent test or a new baseline.

\subsection{MMRDR-UWF and multi-source fine-tuning stress tests}
The separate MMRDR-UWF local development subset contains 1,172 images (572 referable and 600 nonreferable at the grade-$\geq2$ boundary) from the multimodal MMRDR release \cite{mmrdr}. Patient identifiers were not available in the exported rows, so the local split is image-level; the official test partition was not evaluated. Candidate and threshold comparisons on this subset are therefore development-only. A later four-source fine-tuning run used DDR, DeepDRiD, DRTiD, and IDRiD training images \cite{idrid}, with a 1,337-image-row manifest spanning 982 eyes. A subsequent audit found that all 981 historical reference eyes appeared in its validation rows. This correction supersedes an earlier claim that the reference eyes were held apart; the fine-tuning experiment is exposed relative to the historical reference and cannot provide independent validation.

\subsection{Statistics}
Specificity is reported as an empirical ratio with its numerator and denominator. We do not attach a population confidence interval to the headline result because threshold and reference selection reuse the cohort and patient dependence cannot be ignored. A standard exact binomial interval describes sampling for a fixed procedure but does not correct post-selection bias in this setting \cite{clopperpearson,cawley2010}. The 75\% specificity target corresponds arithmetically to at least $\lceil0.75\times385\rceil=289$ true negatives on this denominator; it is a development arithmetic target, not a clinical threshold.

\section{Results}
\subsection{The replay is exact, but it is optimistic by construction}
The deterministic replay reproduced the frozen counts: 270 true negatives among 385 gradable nonreferable cases (70.13\%), no observed referable clearances among 463 gradable referable cases, no ungradable case cleared, and 100\% binary action assignment. By source, specificity was 134/206 (65.05\%) for DDR, 51/52 (98.08\%) for DeepDRiD, and 85/127 (66.93\%) for DRTiD. The large source spread and unequal source composition make the pooled estimate sensitive to case mix. More importantly, the score reference vectors and thresholds were estimated on the same cases, so the point estimate is a resubstitution-style development result.

\subsection{Patient-grouped procedures do not retain the zero-miss headline}
\begin{table}[t]
\centering
\caption{Specificity and observed misses across development procedures on the inspected 981-case cohort. Patient grouping improves the evaluation design but does not make this cohort untouched.}
\label{tab:procedures}
\begin{tabular}{lrrrl}
\toprule
Procedure & TN/385 & Sp. & Misses & Status \\
\midrule
Frozen same-cohort replay & 270 & 70.13\% & 0 & exposed replay \\
Fixed blend, 5-fold grouped & 268 & 69.61\% & 3 & development CV \\
Nested weights, 5$\times$4 & 248 & 64.42\% & 5 & development CV \\
Nested weights + margin & 253 & 65.71\% & 6 & development CV \\
Post-hoc zero-miss margin & 243 & 63.12\% & 0 & held-out tuned \\
\bottomrule
\end{tabular}
\end{table}

The fixed-blend patient-grouped estimate is 0.52 percentage points below the same-cohort replay and includes three missed referable-or-ungradable cases. Nested weight selection decreases specificity by 5.71 points from the replay and yields five misses. Adding the nested margin increases specificity from 64.42\% to 65.71\% but also increases misses to six. The post-hoc margin sweep can force observed misses to zero, but specificity falls to 63.12\%; because the margins were selected on the held-out folds, even that zero-miss point is optimistic and cannot be deployed. This is the practical consequence of selection bias: outer-fold labels must not be used to choose the operating point whose outer-fold performance is then reported \cite{cawley2010,varma2006}.

\begin{figure}[t]
\centering
\small
\setlength{\tabcolsep}{3pt}
\begin{tabular}{lrl}
DDR & \textcolor{blue}{\rule{39.0mm}{2.4mm}} & 65.05\% (134/206)\\[2pt]
DeepDRiD & \textcolor{blue}{\rule{58.8mm}{2.4mm}} & 98.08\% (51/52)\\[2pt]
DRTiD & \textcolor{blue}{\rule{40.2mm}{2.4mm}} & 66.93\% (85/127)\\[2pt]
Pooled & \textcolor{black}{\rule{42.1mm}{2.4mm}} & 70.13\% (270/385)
\end{tabular}
\caption{Observed same-cohort specificity by source. Bar lengths are proportional to the exact source ratios in Table~\ref{tab:cohort}; pooled specificity is a case-weighted aggregate over the three source denominators. This figure visualizes exposed development results and is not an external-performance estimate.}
\label{fig:source}
\end{figure}

\subsection{The later candidate does not provide a safe lift}
The clean-development-trained ResNet-50 candidate alone cleared 54/385 nonreferable cases (14.03\%) and missed two referable cases. OR-combining candidate CLEAR actions with the frozen baseline increased true negatives from 270 to 273 (70.91\%, a 0.78-point descriptive increase) but still missed two referable cases. The 75\% arithmetic target would require 289 true negatives; the union remains 16 short and violates the zero-miss condition. The baseline remains frozen and unchanged.

\subsection{The apparent gap to 75\% is not a license to tune on labels}
On the frozen denominator, 75\% requires 19 additional true negatives beyond the original 270. Later exposed-cohort experiments reduce this arithmetic gap but either introduce misses or use thresholds selected after inspecting reference labels. No result in this study establishes a 75\% zero-miss candidate on new cases.

\subsection{A second domain stress test does not remove the evidence gap}
On the MMRDR-UWF validation subset, the best tested RETFound feature head at a within-validation zero-miss operating point cleared 63/600 negatives (10.50\%) with AUROC 0.919; an RBF-SVM cleared 60/600 (10.00\%). The internal ResNet+DINO/448-pixel ResNet reference cleared 202/600 (33.67\%; AUROC 0.932). Adding 25\% of the RETFound-logistic rank produced 207/600 (34.50\%; AUROC 0.940), five additional clearances. The representation was derived from the retinal foundation model RETFound \cite{retfound}. Since the same validation labels informed both candidate and threshold comparison, these are internal development estimates, not a held-out model comparison.

On the audited 982-eye fine-tuning validation manifest, the MMRDR-initialized warm start reached 95/386 gradable negatives at zero observed misses (24.61\%). Fine-tuning at epochs 1, 2, and 3 yielded 63/386 (16.32\%), 74/386 (19.17\%), and 73/386 (18.91\%), respectively. The warm start remained the strongest tested point, but every historical reference eye was in the validation rows. These figures are reported to preserve the provenance correction and training outcome; they do not establish an independent gain or a clean epoch-selection comparison.

\section{Discussion}
The key finding is methodological: a zero-miss specificity point can look stable in replay and then lose zero-miss status when calibration is moved into patient-grouped folds. The observed 70.13\% should therefore be retained as a historical frozen reference, not promoted to an independent performance claim. Patient grouping is necessary to prevent same-patient leakage across folds, but it cannot undo earlier cohort-wide exploration. This distinction is often obscured when a dataset split is called ``test'' even after repeated model repair.

The source-specific results also argue against interpreting the pooled metric as a universal operating point. DeepDRiD contributes a small negative denominator (52 cases) with 51 correctly cleared, while DDR and DRTiD provide much larger denominators and lower specificities. A pooled value summarizes this particular composition; it does not guarantee equivalent behavior in a new clinic, camera, geography, or prevalence setting.

The later transfer experiment is a useful negative result. A candidate trained with patient groups excluded from training and threshold selection can add a few clearances on an exposed cohort, yet the simple OR policy creates two misses. In the original replay, 270 of 385 nonreferable cases are cleared and 115 are referred; the 463 gradable referable and 133 ungradable records are referred. Small specificity gains are not acceptable when they break the safety constraint. A clinically meaningful next step is not another search over this reference pool; it is to freeze one complete image-to-action artifact and evaluate once on an untouched patient-level external cohort.

The MMRDR experiments extend this caution across a second acquisition domain, but they are not evidence that a particular architecture or dataset is generally ineffective. Within-validation zero-miss points were modest, and the MMRDR-initialized training results were affected by a validation manifest that included the historical reference eyes. Their value here is diagnostic: higher AUROC, a few extra clearances, or additional training data cannot substitute for a frozen procedure evaluated on a genuinely untouched cohort.

\subsection{What this study does and does not establish}
It establishes deterministic replay of a source-conditional score policy and documents the difference between resubstitution-style and patient-grouped development procedures on this historical cohort. It does not establish a calibrated probability, transportable zero-miss behavior, a regulatory claim, clinical utility, prospective workflow safety, or a complete raw-image model. The word ``coverage'' is used only for action assignment when specified; it must not be confused with the fraction of cases autonomously diagnosed by a validated system.

\section{Limitations}
First, the headline policy was built and tuned on the same 981-case cohort used for replay. Second, the source component scores are precomputed; this artifact is not end-to-end inference, and all source checkpoints/preprocessing contracts are not available together in this checkout. Third, label definitions and image acquisition differ across datasets, and the pooled result is not prevalence- or population-standardized. Fourth, patient-grouped folds only protect their own fitting operations; prior project-wide inspection remains. Fifth, the current source-level negatives are modest, especially for DeepDRiD. Sixth, the transfer union is a diagnostic construction, not an adopted policy. Finally, no clinical trial, prospective deployment, or independent external certification is included.

\section{Reproducibility and Data Availability}
The audit script has two modes. Its clean-clone default, \path{--mode summary}, checks selected headline counts in tracked summary files and hashes those files; it does not execute the policy or independently reconstruct the counts from patient records. The optional \path{--mode full-replay} uses \path{results/deployable_exposed_baseline_7013_20261013_with_components/policy.json} and \path{validation_replay.csv}, with the evaluator in \path{src/active_retina/exposed_expert_policy.py}, to recompute every action and risk score for the 981 supplied component-score rows and check source counts and saved summaries. Grouped, nested, and transfer summaries are listed in the reproduction-package README. Both modes use Python 3.12 and the dependency pinned in \path{paper/reproducibility/requirements.txt}; the full replay also writes SHA-256 digests for its inputs. Neither mode recreates image-level component scores, retrains models, or establishes external validation.

The clean public checkout omits the local policy and per-eye score/label rows pending author confirmation that redistribution is permitted by every applicable source term; consequently, the public aggregate check is not a clean-clone reproduction of the row-level replay. The DDR repository states CC BY-NC-SA 4.0 for dataset materials; the DeepDRiD code repository carries CC BY-SA 4.0; and DRTiD is described by its official repository as research-only and application-accessed. Those statements do not, by themselves, settle the rights for this project's derived per-eye outputs. MMRDR and IDRiD item-level terms also require author verification. Raw images are not redistributed. The code-and-data release checklist is \path{paper/SUBMISSION_RELEASE_CHECKLIST.md}; until its rights and author declarations are resolved, the full reproduction bundle and final author/data-use checks remain incomplete. No independent external validation cohort is reported.

\section{Conclusion}
The 70.13\% specificity result is exactly reproducible as a frozen score-level replay on its 981-case development cohort, but that same cohort supplied its percentile references and operating thresholds. Patient-grouped procedures yield 64.42--69.61\% specificity with observed misses; a post-hoc zero-miss frontier falls to 63.12\% after held-out tuning. A later candidate-transfer union reaches 70.91\% on the same exposed cohort but introduces two misses. MMRDR-UWF and fine-tuning stress tests also remain internal development evidence, with explicit image-level split and manifest-overlap limitations. The defensible conclusion is not that ActiveRetina has a clinically ready 70.13\% model. It is that retrospective zero-miss operating points require strict provenance and a new untouched evaluation before they can support transportable claims.

\appendix
\section{Compact Claim Register}
\begin{table}[h]
\centering
\caption{Claim-to-evidence boundary.}
\begin{tabularx}{\linewidth}{lX}
\toprule
Claim & Status and evidence boundary \\
\midrule
70.13\% replay & Reproduced on the same exposed 981-case pool; not independent validation.\\
Zero observed misses & True for that replay; thresholds were selected using the same cohort.\\
Patient-generalized 70.13\% & Not established; grouped procedures differ and have observed misses.\\
75\% zero-miss performance & Not achieved; no candidate meets both conditions on an untouched cohort.\\
Clinical autonomy & Not established; score-level evaluator is not full image-to-output deployment.\\
\bottomrule
\end{tabularx}
\end{table}

\section{Submission Readiness Note}
This is a coherent research-report draft with explicit methods, results, negative findings, and a score-level replay auditor. It is not yet an arXiv release package: the clean-clone reproduction inputs are not included pending a documented rights decision for derived per-eye scores and labels, and author/data-use declarations remain to be confirmed. The release checklist records the exact remaining items. Before posting, authors must verify reference metadata, dataset terms, contribution and affiliation details, disclosures, the repository revision, and correspondence between source and compiled PDF. The exposed-cohort boundary must remain in the abstract and conclusion. The current evidence does not justify clinical deployment claims.

\begin{thebibliography}{20}
\bibitem{gulshan2016} V. Gulshan et al., ``Development and validation of a deep learning algorithm for detection of diabetic retinopathy in retinal fundus photographs,'' \emph{JAMA}, vol. 316, no. 22, pp. 2402--2410, 2016.
\bibitem{ting2017} D. S. W. Ting et al., ``Development and validation of a deep learning system for diabetic retinopathy and related eye diseases using retinal images from multiethnic populations with diabetes,'' \emph{JAMA}, vol. 318, no. 22, pp. 2211--2223, 2017.
\bibitem{abramoff2018} M. D. Abr\`amoff et al., ``Pivotal trial of an autonomous AI-based diagnostic system for detection of diabetic retinopathy in primary care offices,'' \emph{npj Digital Medicine}, vol. 1, art. no. 39, 2018.
\bibitem{bellemo2019} V. Bellemo et al., ``Artificial intelligence using deep learning to screen for referable and vision-threatening diabetic retinopathy in Africa,'' \emph{The Lancet Digital Health}, vol. 1, no. 1, pp. e35--e44, 2019.
\bibitem{deepdrid} Y. Liu et al., ``DeepDRiD: Diabetic retinopathy-grading and image quality estimation challenge,'' \emph{Patterns}, vol. 3, no. 6, art. no. 100512, 2022, doi: 10.1016/j.patter.2022.100512.
\bibitem{ddr} T. Li et al., ``Diagnostic assessment of deep learning algorithms for diabetic retinopathy screening,'' \emph{Information Sciences}, vol. 501, pp. 511--522, 2019.
\bibitem{geifman2017} Y. Geifman and R. El-Yaniv, ``Selective classification for deep neural networks,'' in \emph{Advances in Neural Information Processing Systems}, 2017.
\bibitem{selectivenet} Y. Geifman and R. El-Yaniv, ``SelectiveNet: A deep neural network with an integrated reject option,'' in \emph{Proceedings of the 36th International Conference on Machine Learning}, 2019, pp. 2151--2159.
\bibitem{vovk2005} V. Vovk, A. Gammerman, and G. Shafer, \emph{Algorithmic Learning in a Random World}. Springer, 2005.
\bibitem{angelopoulos2021} A. N. Angelopoulos and S. Bates, ``A gentle introduction to conformal prediction and distribution-free uncertainty quantification,'' arXiv:2107.07511, 2021.
\bibitem{angelopoulos2024} A. N. Angelopoulos, S. Bates, E. J. Cand\`es, M. I. Jordan, and L. Lei, ``Learn then test: Calibrating predictive algorithms to achieve risk control,'' \emph{The Annals of Applied Statistics}, vol. 19, no. 2, pp. 1641--1662, 2025, doi: 10.1214/24-AOAS1998.
\bibitem{tripodai} G. S. Collins et al., ``TRIPOD+AI statement: updated guidance for reporting clinical prediction models that use regression or machine learning methods,'' \emph{BMJ}, vol. 385, q902, 2024.
\bibitem{claim} B. J. Mongan, L. Moy, and C. E. Kahn, ``Checklist for Artificial Intelligence in Medical Imaging (CLAIM): a guide for authors and reviewers,'' \emph{Radiology: Artificial Intelligence}, vol. 2, no. 2, e200029, 2020.
\bibitem{clopperpearson} C. J. Clopper and E. S. Pearson, ``The use of confidence or fiducial limits illustrated in the case of the binomial,'' \emph{Biometrika}, vol. 26, no. 4, pp. 404--413, 1934.
\bibitem{cawley2010} G. C. Cawley and N. L. C. Talbot, ``On over-fitting in model selection and subsequent selection bias in performance evaluation,'' \emph{Journal of Machine Learning Research}, vol. 11, pp. 2079--2107, 2010.
\bibitem{varma2006} S. Varma and R. Simon, ``Bias in error estimation when using cross-validation for model selection,'' \emph{BMC Bioinformatics}, vol. 7, art. no. 91, 2006, doi: 10.1186/1471-2105-7-91.
\bibitem{roberts2017} D. R. Roberts et al., ``Cross-validation strategies for data with temporal, spatial, hierarchical, or phylogenetic structure,'' \emph{Ecography}, vol. 40, no. 8, pp. 913--929, 2017, doi: 10.1111/ecog.02881.
\bibitem{drtidpaper} ``Cross-field transformer for diabetic retinopathy grading on two-field fundus images,'' arXiv:2211.14552, 2022.
\bibitem{idrid} P. Porwal et al., ``Indian Diabetic Retinopathy Image Dataset (IDRiD): A database for diabetic retinopathy screening research,'' \emph{Data}, vol. 3, no. 3, art. no. 25, 2018, doi: 10.3390/data3030025.
\bibitem{mmrdr} Z. Tang et al., ``A multimodal retinal image dataset for diabetic retinopathy detection using foundation models,'' \emph{Scientific Data}, vol. 13, no. 1, art. no. 639, 2026, doi: 10.1038/s41597-026-07005-9.
\bibitem{retfound} Y. Zhou et al., ``A foundation model for generalizable disease detection from retinal images,'' \emph{Nature}, vol. 622, pp. 156--163, 2023, doi: 10.1038/s41586-023-06555-x.
\bibitem{galappaththige2024} C. J. Galappaththige, G. Kuruppu, and M. H. Khan, ``Generalizing to unseen domains in diabetic retinopathy classification,'' in \emph{Proc. IEEE/CVF Winter Conf. Applications of Computer Vision}, 2024, pp. 7685--7695.
\bibitem{claim2024} A. S. Tejani et al., ``Checklist for Artificial Intelligence in Medical Imaging (CLAIM): 2024 update,'' \emph{Radiology: Artificial Intelligence}, vol. 6, no. 4, e240300, 2024, doi: 10.1148/ryai.240300.
\end{thebibliography}
\end{document}
